\section{Formulaci\'on matem\'atica del problema geneal\'ogico del continuo}
\label{sec:formulacion-genealogica-continuo}

\subsection{Magnitud, completitud, cardinalidad y g\'enesis}

La formulaci\'on matem\'atica del continuo comprende problemas de naturaleza
distinta que deben conservarse separados. La teor\'ia de las proporciones
determina relaciones exactas entre magnitudes; la completitud de
\(\mathbb R\) garantiza la existencia de l\'imites de sucesiones de Cauchy y
de familias encajadas de intervalos; la topolog\'ia caracteriza conexi\'on,
compacidad y continuidad; la teor\'ia de la medida asigna contenido a
subconjuntos; y la teor\'ia de conjuntos compara cardinales. Ninguna de estas
estructuras, considerada aisladamente, especifica la ley de formaci\'on que
produce una coordenada real a partir de configuraciones finitas y conserva
simult\'aneamente su historia de construcci\'on.

La hip\'otesis cl\'asica del continuo es el enunciado cardinal
\begin{equation}
 2^{\aleph_0}=\aleph_1,
 \label{eq:CH-clasica}
\end{equation}
equivalente, en ZFC, a que todo subconjunto no numerable de
\(\mathbb R\) posea cardinal \(2^{\aleph_0}\). Los teoremas de G\"odel y
Cohen establecen, bajo la hip\'otesis metamatem\'atica de consistencia de ZFC,
que \eqref{eq:CH-clasica} y su negaci\'on poseen modelos de ZFC. Esta
independencia concierne al cardinal del conjunto de todas las partes de
\(\mathbb N\). El problema geneal\'ogico adopta una variable diferente:
determina un sistema de configuraciones, prolongaciones y lecturas del cual
\(\mathbb R\) resulta como cociente topol\'ogico. La distinci\'on entre ambos
problemas es formal: uno compara cardinales en un universo conjuntista; el
otro construye el portador, sus fibras de historia y su evaluaci\'on.

\begin{definition}[Problema geneal\'ogico del continuo]
Una resoluci\'on geneal\'ogica del continuo consta de los siguientes datos:
\begin{enumerate}[label=\textup{(\roman*)}]
 \item una familia de espacios finitos de configuraciones
       \(\{\mathcal E_k\}_{k\geq0}\) y aplicaciones de restricci\'on
       \(\rho_{k+1,k}:\mathcal E_{k+1}\to\mathcal E_k\);
 \item una ley de prolongaci\'on que determine qu\'e configuraciones de
       profundidad \(k+1\) conservan el estado de profundidad \(k\);
 \item un estado enriquecido que retenga los invariantes necesarios para
       distinguir configuraciones con igual coordenada visible;
 \item una familia compatible de cilindros de di\'ametro convergente a cero;
 \item una evaluaci\'on arquimediana sobre el l\'imite inverso y una
       equivalencia que identifique exactamente las historias con igual
       valor;
 \item funtores que transporten, sobre la misma historia, estructura
       espectral, balance, incidencia, cohomolog\'ia y determinante;
 \item una extensi\'on a todas las cartas compactas de la recta y a los
       rangos dimensionales obtenidos por cocientes de memoria.
\end{enumerate}
\end{definition}

La Holograf\'ia Modular Tri\'adica realiza esos datos mediante la cadena
tipada
\begin{equation}
 \mathcal D_9
 \longrightarrow \APP
 \longrightarrow \TRIT
 \longrightarrow \TPK
 \longrightarrow X_{\TPK}^{\rm enr}
 \longrightarrow \Cdisc,
 \qquad \mathcal D_9=\{1,\ldots,9\}.
 \label{eq:cadena-rectora-formulacion}
\end{equation}
El alfabeto \(\mathcal D_9\) proporciona representantes positivos de las
clases non\'adicas. APP construye el soporte bidimensional y sus dos hojas
aritm\'eticas; TRIT orienta la transici\'on local y conserva su acarreo
ternario; TPK compone selecci\'on, transporte, actualizaci\'on, retorno y
prolongaci\'on sobre el estado enriquecido. El objeto \(\Cdisc\) es el
l\'imite conjunto de las configuraciones compatibles. La evaluaci\'on real y
la realizaci\'on excepcional act\'uan despu\'es sobre ese mismo antecedente.

\subsection{Unidad, extensi\'on y conservaci\'on geneal\'ogica}

Sea \(P_9\) el camino orientado de nueve fases y sea \(C_9\) su cierre. Una
vuelta de \(C_9\) restituye la coordenada de fase y aumenta el contador de
vueltas. Esta separaci\'on elemental induce dos invariantes:
\begin{equation}
 r(n)=1+((n-1)\bmod9),
 \qquad
 q(n)=\left\lfloor\frac{n-1}{9}\right\rfloor,
 \qquad
 n=9q(n)+r(n).
 \label{eq:division-nonadica-inaugural}
\end{equation}
La identidad
\begin{equation}
 r(n+9)=r(n),
 \qquad q(n+9)=q(n)+1
 \label{eq:retorno-inaugural}
\end{equation}
expresa una periodicidad de la primera coordenada y una traslaci\'on de la
segunda. La conservaci\'on de la unidad a trav\'es del refinamiento se
formaliza mediante pares secci\'on--retracci\'on. Para
\(\mathcal C_9^{(d)}=(\mathbb Z/9\mathbb Z)^d\),
\begin{equation}
 h_{d,\star}(x)=(x,\star),
 \qquad
 \pi_{d+1,d}(x_1,\ldots,x_{d+1})=(x_1,\ldots,x_d),
 \qquad
 \pi_{d+1,d}\circ h_{d,\star}=\id.
 \label{eq:conservacion-dimensional-unidad}
\end{equation}
Cada nivel superior contiene una secci\'on isomorfa al nivel precedente y
una coordenada adicional de memoria. La proyecci\'on uniforme satisface
\((\pi_{d+1,d})_*\mu_{d+1}=\mu_d\). La extensi\'on dimensional conserva,
por tanto, estructura y medida antes de introducir una realizaci\'on
m\'etrica.

\begin{definition}[Simetr\'ia geneal\'ogica]
Una simetr\'ia geneal\'ogica es una familia de biyecciones
\(g_k:\mathcal E_k\to\mathcal E_k\) que preserva los tipos del estado y
satisface
\begin{equation}
 \rho_{k+1,k}\circ g_{k+1}=g_k\circ\rho_{k+1,k}.
 \label{eq:simetria-genealogica-formulacion}
\end{equation}
La familia induce un automorfismo del l\'imite inverso. Si, adem\'as,
preserva la medida de cilindros y los funtores de salida, conserva la unidad
en todas sus realizaciones compatibles.
\end{definition}

Esta definici\'on integra las nociones de escala y dimensi\'on en la
estructura de prolongaci\'on. Una escala especifica la profundidad y el
di\'ametro del cilindro evaluado. Una dimensi\'on especifica el rango de
las memorias independientes que sobreviven a las relaciones de frontera. La
variaci\'on de escala refina una coordenada; la elevaci\'on dimensional a\~nade
un grado de libertad recuperable. Ambas operaciones quedan determinadas por
morfismos distintos.

\subsection{Teorema rector de construcci\'on}

\begin{theorem}[Construcci\'on geneal\'ogica del continuo]
\label{thm:rector-construccion-continuo}
Para cada \(m\geq1\), sea
\[
 \mathcal E_0^{(m)}
 \xleftarrow{\rho_{1,0}}
 \mathcal E_1^{(m)}
 \xleftarrow{\rho_{2,1}}\cdots
\]
un sistema inverso de espacios finitos no vac\'ios con aplicaciones
sobreyectivas. Sup\'ongase que cada \(x_k\in\mathcal E_k^{(m)}\) determina un
intervalo compacto \(I_k(x_k)\subseteq K_m=[-m,m]\) y que existen
\(\delta_k\downarrow0\) tales que
\begin{align}
 \rho_{k+1,k}(x_{k+1})=x_k
 &\Longrightarrow I_{k+1}(x_{k+1})\subseteq I_k(x_k),
 \label{eq:encaje-cilindros-rector}\\
 \operatorname{diam}I_k(x_k)&\leq\delta_k,
 \label{eq:contraccion-cilindros-rector}\\
 I_k(x_k)&=
 \bigcup_{\rho_{k+1,k}(y)=x_k}I_{k+1}(y),
 \qquad
 K_m=\bigcup_{x_0\in\mathcal E_0^{(m)}}I_0(x_0).
 \label{eq:cobertura-cilindros-rector}
\end{align}
\Needspace{8\baselineskip}
Entonces:
\begin{enumerate}[label=\textup{(\alph*)}]
 \item el l\'imite
 \(\Omega_m=\varprojlim_k\mathcal E_k^{(m)}\) es compacto y no vac\'io;
 \item la aplicaci\'on
 \begin{equation}
  \operatorname{Val}_m((x_k)_k)=
  \text{el elemento de }\bigcap_{k\geq0}I_k(x_k)
  \label{eq:valor-rector}
 \end{equation}
 es continua y sobreyectiva sobre \(K_m\);
 \item la relaci\'on
 \(x\sim_{\mathbb R}y\Longleftrightarrow
 \operatorname{Val}_m(x)=\operatorname{Val}_m(y)\) induce un homeomorfismo
 \begin{equation}
  \Omega_m/\!\sim_{\mathbb R}\;\cong K_m;
  \label{eq:cociente-rector-compacto}
 \end{equation}
 \item la compatibilidad de las inclusiones \(K_m\hookrightarrow K_{m+1}\)
 produce
 \begin{equation}
  \varinjlim_m\bigl(\Omega_m/\!\sim_{\mathbb R}\bigr)
  \cong\mathbb R.
  \label{eq:agotamiento-real-rector}
 \end{equation}
\end{enumerate}
Si sobre cada \(\mathcal E_k^{(m)}\) act\'ua un ensamblaje natural
\(\mathcal O_k\) de transporte polar, balance firmado, incidencia celular,
complejo de puertos e invariantes homol\'ogico--determinantales, entonces
existe una aplicaci\'on l\'imite \(\mathcal O_\infty\) sobre \(\Omega_m\).
Una evaluaci\'on arquimediana y un lector excepcional calibrado definen, por
consiguiente, la aplicaci\'on conjunta
\begin{equation}
 \mathfrak D_{\chi}(x)=
 \bigl(\operatorname{Val}_m(x),\mathfrak E_\chi(x)\bigr),
 \qquad
 \mathfrak D_\chi:\Omega_m^{\rm cal}
 \longrightarrow K_m\times\mathbf{Exc}.
 \label{eq:doble-salida-teorema-rector}
\end{equation}
\end{theorem}

\begin{proof}
Los espacios finitos son compactos y las aplicaciones de enlace son
continuas y sobreyectivas; el teorema del l\'imite inverso proporciona la
compacidad y la no vacuidad de \(\Omega_m\). Las condiciones
\eqref{eq:encaje-cilindros-rector} y
\eqref{eq:contraccion-cilindros-rector} garantizan que la intersecci\'on de
los cilindros de una secci\'on compatible contiene un \'unico punto. La
cobertura coherente \eqref{eq:cobertura-cilindros-rector} construye una
secci\'on sobre cada punto de \(K_m\), y prueba la sobreyectividad. La
aplicaci\'on inducida por el cociente es una biyecci\'on continua desde un
compacto hacia un espacio de Hausdorff; por tanto, es un homeomorfismo. El
agotamiento por los compactos \([-m,m]\) establece
\eqref{eq:agotamiento-real-rector}. La naturalidad
\(u^{\mathcal O}_{\ell,k}\circ\mathcal O_\ell
=\mathcal O_k\circ u^{\mathcal E}_{\ell,k}\) determina
\(\mathcal O_\infty\) por la propiedad universal del l\'imite inverso. Los
dos lectores de \eqref{eq:doble-salida-teorema-rector} poseen el mismo
dominio y, por ello, constituyen componentes de una sola aplicaci\'on.
\end{proof}

El teorema identifica tres niveles matem\'aticos. La secci\'on del l\'imite
inverso es el n\'umero geneal\'ogico; la clase de equivalencia es el valor
arquimediano; y el par formado por valor e incidencia es la doble
realizaci\'on de la misma historia. La recta real es la realizaci\'on
topol\'ogica del cociente, mientras que el sistema proyectivo conserva la
informaci\'on eliminada por ese cociente.

\subsection{Obligaciones de una resoluci\'on completa}

La formulaci\'on anterior fija condiciones verificables. La construcci\'on
debe especificar los espacios finitos, las aplicaciones de enlace, la ley de
prolongaci\'on y el estado conservado; debe demostrar no vacuidad,
compatibilidad y contracci\'on de cilindros; debe identificar el cociente con
la totalidad de \(\mathbb R\), incluidos los valores negativos y los
extremos de las cartas; debe distinguir cardinalidad, ordinalidad, escala y
dimensi\'on; y debe probar la naturalidad de las realizaciones que descienden
del mismo estado. APP, TRIT y TPK proporcionan los objetos requeridos para
estas obligaciones. Las secciones siguientes fijan sus dominios,
codominios y leyes de composici\'on.
